%%=============================================================================
%% Verantwoording van het gebruik van AI - VERPLICHT
%%=============================================================================

\chapter{\IfLanguageName{dutch}{Verantwoording van het gebruik van AI}{Use of AI}}%
\label{ch:ai}

%% Deze bijlage is verplicht, ook als je geen AI gebruikt hebt. Schrijf dan
%% gewoon dat je er geen gebruikt hebt.
%%
%% De opleiding verwacht uitdrukkelijk dat je AI doelgericht inzet om iets
%% nieuws te maken. Het is dus een competentie die je toont, geen bekentenis
%% die je aflegt. Wat wél telt: je moet kunnen uitleggen wat je gebruikt hebt,
%% waarvoor, en wat je zelf gecontroleerd hebt.
%%
%% Reken erop dat de jury hiernaar vraagt. "Wat is hiervan van jou?" is een van
%% de vaste vragen tijdens de verdediging.

Vervang deze alinea door een korte inleiding: heb je AI gebruikt tijdens dit
project, en zo ja, welke rol speelde ze in je werkwijze?

\section{\IfLanguageName{dutch}{Waarvoor en hoe}{What and how}}%
\label{sec:ai-waarvoor}

%% Vul de tabel aan met je eigen gebruik. Eén regel per soort gebruik, niet
%% per gesprek. Wees concreet in de laatste kolom: die is het belangrijkst.

\begin{table}[htbp]
  \centering
  \begin{tabular}{p{0.22\textwidth}p{0.2\textwidth}p{0.44\textwidth}}
    \toprule
    \textbf{Waarvoor} & \textbf{Welk hulpmiddel} & \textbf{Wat ik ermee gedaan heb, en wat ik controleerde} \\
    \midrule
    Voorbeeld: opzet van de gegevenslaag
      & naam en versie
      & Voorstel laten genereren, daarna zelf herschreven omdat de foutafhandeling ontbrak; met tests nagekeken. \\
    Voorbeeld: uitleg over een foutmelding
      & naam en versie
      & Als uitleg gebruikt, oplossing zelf geschreven na controle in de officiële documentatie. \\
    Voorbeeld: nalezen van dit rapport
      & naam en versie
      & Enkel op spelling en zinsbouw; de inhoud en de structuur zijn van mij. \\
    \bottomrule
  \end{tabular}
  \caption[Gebruik van AI in dit project]{\label{tab:ai}%
    Overzicht van het gebruik van AI in dit project.}
\end{table}

\section{\IfLanguageName{dutch}{Wat van mij is}{What is my own work}}%
\label{sec:ai-eigenwerk}

%% Twee regels die de opleiding hanteert, en die je hier bevestigt:
%%
%% 1. Code die je niet kan uitleggen, hoort niet in je oplevering. Je bent
%%    verantwoordelijk voor alles wat in je repository staat, ook voor wat een
%%    hulpmiddel geschreven heeft.
%% 2. De analyse, de keuzes, de afbakening en de tekst van dit rapport zijn van
%%    jou. AI mag je helpen formuleren; ze mag niet in jouw plaats denken.
%%
%% Geef hier ook aan hoe je in je repository zichtbaar maakt welke stukken met
%% AI tot stand kwamen - bijvoorbeeld met een korte notitie in de commit of een
%% opmerking bovenaan het bestand.

\section{\IfLanguageName{dutch}{Zorgvuldigheid en regelgeving}{Diligence and regulation}}%
\label{sec:ai-regelgeving}

%% Beantwoord hier kort de drie vragen die er in de praktijk toe doen.
%%
%% 1. HEB JE BEDRIJFSGEGEVENS IN EEN AI-HULPMIDDEL GEPLAKT?
%%    Broncode, klantgegevens, configuratie of persoonsgegevens ingeven in een
%%    publieke dienst is een doorgifte aan een derde partij. Dat kan in strijd
%%    zijn met je geheimhoudingsplicht en met de AVG (de Algemene Verordening
%%    Gegevensbescherming, ook bekend als de GDPR). Vraag je mentor wat bij
%%    jouw bedrijf toegelaten is, en noteer het antwoord hier.
%%
%% 2. ZIT ER AI IN WAT JE GEBOUWD HEBT?
%%    Dan is de Europese AI-verordening (de AI Act, verordening 2024/1689) van
%%    belang: haar verplichtingen zijn sinds augustus 2026 grotendeels van
%%    toepassing. Ga na in welke risicocategorie je toepassing valt, en of er
%%    transparantieplichten gelden - bijvoorbeeld wanneer gebruikers met een
%%    AI-systeem praten of wanneer je toepassing inhoud genereert. In beide
%%    gevallen moet de gebruiker dat weten.
%%    Gebruik je AI enkel als hulpmiddel bij het programmeren, dan valt je
%%    project daar meestal niet onder; schrijf dat dan gewoon zo op.
%%
%% 3. WAT MET AUTEURSRECHT EN LICENTIES?
%%    Gegenereerde code kan sterk lijken op bestaande code met een licentie.
%%    Vermeld hoe je daarmee omgegaan bent.

Vervang deze alinea door je antwoord op de drie vragen hierboven.

%% Heb je geen AI gebruikt? Vervang dit hele hoofdstuk dan door één zin:
%% "Bij de realisatie van dit project is geen gebruikgemaakt van generatieve AI."
